% Contenido del póster: Miglo (2025), Journal of Economic Studies. DOI 10.1108/jes-03-2024-0194
% Elaborado con apoyo de IA a partir del resumen oficial del artículo (OpenAlex), 2026-11-12.
\def\DOI{10.1108/jes-03-2024-0194}
\def\TituloEs{Crowdfunding, información asimétrica y orden jerárquico de financiamiento}
\def\TituloOriginal{Crowdfunding, asymmetric information and pecking order}
\def\Autores{Miglo}
\def\Anio{2025}
\def\Revista{Journal of Economic Studies}

\def\Problema{Existe una brecha entre la teoría y la evidencia sobre el crowdfunding: los volúmenes son pequeños frente al financiamiento tradicional y llegan muchas empresas de baja calidad, algo que las teorías previas no explican del todo.}
\def\Objetivo{Proponer una teoría de cómo elige el emprendedor entre deuda, capital y crowdfunding de recompensa cuando la información sobre la calidad de sus productos es asimétrica.}
\def\Metodo{Modelo teórico de selección de fuentes de financiamiento con información asimétrica sobre la calidad esperada del producto. Los pagos del crowdfunding de recompensa son preventas que se reciben antes de invertir en el proyecto. El enfoque se compara con modelos en los que la asimetría es sobre el costo de producción.}
\def\Resultados{\begin{itemize}
  \item En la mayoría de los casos, el crowdfunding queda al final del orden jerárquico de financiamiento.
  \item Como las preventas llegan antes de la inversión, la ganancia del emprendedor no depende de la calidad real del producto.
  \item Eso genera subsidios cruzados de los emprendedores de alta calidad hacia los de baja calidad en un equilibrio agrupador.
  \item El modelo es coherente con los volúmenes pequeños del crowdfunding y con la entrada excesiva de empresas de baja calidad.
\end{itemize}}
\def\Conclusion{La información asimétrica explica por qué el crowdfunding de recompensa suele ser la última opción de financiamiento y atrae a proyectos de menor calidad.}
\def\Relevancia{El modelo trata el crowdfunding de recompensa, que en el Perú no forma parte del financiamiento participativo financiero regulado por la SMV. Aun así, su idea central, que la información asimétrica limita el volumen del crowdfunding frente al crédito tradicional, sirve de marco para interpretar el tamaño de un mercado peruano que financió S/~\cifra{monto2025} millones en 2025.}
\def\PalabrasClave{crowdfunding, información asimétrica, orden jerárquico, selección adversa, financiamiento empresarial}
% --- Ficha del estudio (herramientas/posteres/ficha_articulos.py) ---
\def\Ambito{Teórico}
\def\Enfoque{Modelo de información asimétrica (crowdfunding de recompensa)}
\def\FiguraPeru{1}
% --- fin de la ficha ---
